\documentclass[12pt]{article}
\usepackage[utf8]{inputenc}
\usepackage[T1]{fontenc}
\usepackage{hyperref}
\title{Pisidian Language Resources for K-12 Curriculum Development}
\author{Compiled by Claude Sonnet 5.5 (claude-sonnet-5-5) for Andr\'as Kornai}
\date{}
\begin{document}
\maketitle

\section{Introduction}
Pisidian (ISO 639-3 xps; Glottocode pisi1234) is an extinct Anatolian language of the Indo-European family. It is poorly attested. Only about two dozen short inscriptions are known, mostly tomb texts from Sofular (ancient Tymbrias). Others come from Selge, Kesme and Değirmenözü. The texts are largely genealogical, consisting of strings of names, and they are written in a Greek-like alphabet. Sources disagree on the date of the inscriptions. Britannica gives the 1st and 2nd centuries BCE, while a Leiden source gives the 1st and 2nd centuries CE. A comprehensive curriculum conducted in Pisidian is not possible. The surviving record is too small.

\section{Primary Educational Resources}
No educational materials, primers or school programmes for Pisidian were found.

\section{Linguistic and Reference Materials}
\begin{itemize}
\item Wikipedia overview: \url{https://en.wikipedia.org/wiki/Pisidian_language}
\item Britannica entry: \url{https://www.britannica.com/print/article/1380518}
\item Leiden scholarly publication: \url{https://scholarlypublications.universiteitleiden.nl/access/item%3A3677187/view}
\item G. E. Bean, ``Notes and Inscriptions from Pisidia, Part 1'', Anatolian Studies: \url{https://www.cambridge.org/core/journals/anatolian-studies/article/abs/notes-and-inscriptions-from-pisidia-part-1/39C43B2A31BCBDD3B4D24CD125E71674}
\item Unpublished inscriptions from Pisidian Antioch, Anatolian Studies: \url{https://resolve.cambridge.org/core/journals/anatolian-studies/article/abs/unpublished-inscriptions-from-pisidian-antioch/07CF22EBABA6A5CD6B9261121BADF45A}
\end{itemize}
The search results also mention Robert Shafer's 1950 article ``Pisidian'' in the American Journal of Philology and a 2025 Brill chapter by Merlin and Pisaniello. No URLs were found for either.

\section{Religious and Community Resources}
None found.

\section{Digital and Audio Resources}
None found.

\section{Curriculum Development Context}
Pisidian could be used in a school only as a small unit on ancient Anatolian languages. That unit would sit alongside Lycian and Sidetic and would be taught from the epigraphic evidence. There is no basis for a full language curriculum.
\end{document}